\subsection{替换的准则}

形式数学仅包含明确书写的组装。然而，即使使用缩写符号，严格依照此原则发展数学也会导致极其冗长的推理链条。因此，我们将建立与不定组装相关的准则；每条准则将一次性描述对这些组装进行特定操作序列的最终结果。这些准则因此并非理论所必需；其合理性属于元数学范畴。

元数学本身的发展，在实践中需要使用一些缩写符号，其中部分符号已在前面提及。这些符号中的大多数在数学中也有使用。

我们将使用以下准则，称为代入准则：

CS1. 设A和B为组装，x和 \(x'\) 为字母。若 \(x'\) 不出现在 A 中，则 \((B|x)A\) 与 \((B|x')(x'|x)A\) 相同。

CS2. 设 \(A, B\) 和 \(C\) 为组装，\(x\) 和 \(y\) 为不同的字母\footnote{按照第 1 号所述，“x 和 y 是不同的字母”这一说法是一种语言上的滥用：它表示在所考虑的组装中，x 和 y 表示不同的字母。}。若 \(y\) 不出现在 \(B\) 中，则 \((B|x)(C|y)A\) 与下式相同：

\[
\left(\boldsymbol {C} ^ {\prime} | \boldsymbol {y}\right) (\boldsymbol {B} | \boldsymbol {x}) \boldsymbol {A},
\]

其中 \(C'\) 是组装 \((B|x)C\) 。

CS3. 设 A 为一个组装，x 和 \(x'\) 为字母。若 \(x'\) 不出现在 A 中，则 \(\tau_{\mathbf{x}}(A)\) 与 \(\tau_{\mathbf{x}'}(A')\) 相同，其中 \(A'\) 是组装 \((x'|x)A\) 。

CS4. 设 A 和 B 为组装，x 和 y 为不同的字母。若 x 不出现在 B 中，则 \((B|y)\tau_{x}A\) 与 \(\tau_{x}(A')\) 相同，其中 \(A'\) 是组装 \((B|y)A\) 。

CS5. 设 \(A, B, C\) 为组装，\(x\) 为一个字母。组装 \((C|x) (\neg A), (C|x) (\vee AB), (C|x) (\Longrightarrow AB), (C|x) (sAB)\) （其中 \(s\) 是一个特定符号）分别与 \(\neg A', \vee A'B', \Longrightarrow A'B', sA'B',\) 相同，其中 \(A', B'\) 分别是 \((C|x) A, (C|x) B\) 。

下面以CS2为例说明验证的原理。比较把A变为 \((B|x)(C|y)A\) 的操作与把A变为 \((C'|y)(B|x)A\) 的操作。在这两种操作中，A内不同于x和y的符号都不改变。在A中每一处x出现的位置，两种操作都以B替换x；对于第一种操作这是显然的，对于第二种操作则由y不出现在B中而知。最后，在A中每一处y出现的位置，第一种操作先以C替换y，再在C中每一处x出现的位置以B替换x；这显然等同于在A中每逢y出现便以组装 \((B|x)C\) 替换。

